\subsection{PARTIAL PRODUCTS}

Let \((\mathbf{X}_i)_{i\in I}\) be a family of sets, and let \(J\) be a subset of \(I\) . The product \(\prod_{i\in J} X_i\) is called a partial product of \(\prod_{i\in I} X_i\) . If \(f\) is a function whose graph \(F\) is a member of \(\prod_{i\in I} X_i\) , then \(F \circ \Delta_J\) (where \(\Delta_J\) is the diagonal of \(J \times J\) ) is the graph of the restriction of \(f\) to \(J\) . Clearly \(F \circ \Delta_J \in \prod_{i\in J} X_i\) ; the mapping \(F \to F \circ \Delta_J\) of \(\prod_{i\in I} X_i\) into \(\prod_{i\in I} X_i\) is called the projection of index \(J\) and is denoted by \(\operatorname{pr}_J\) .

PROPOSITION 5. Let \((\mathbf{X}_{\iota})_{\iota \in \mathbf{I}}\) be a family of sets and let \(\mathbf{J}\) be a subset of \(\mathbf{I}\) . If for each \(\iota \in \mathbf{I}\) we have \(\mathbf{X}_{\iota} \neq \emptyset\) , the projection \(\operatorname{pr}_{\mathbf{J}}\) is a mapping of \(\prod_{\iota \in \mathbf{I}} \mathbf{X}_{\iota}\) onto \(\prod_{\iota \in \mathbf{J}} \mathbf{X}_{\iota}\) .

In view of the remarks made above, it is enough to prove the following proposition :

PROPOSITION 6. Let \((\mathbf{X}_{\iota})_{\iota \in \mathbf{I}}\) be a family of sets such that \(\mathbf{X}_{\iota} \neq \emptyset\) for all \(\iota \in \mathbf{I}\) . If \(g\) is a mapping of \(\mathbf{J} \subset \mathbf{I}\) into \(\mathbf{A} = \bigcup_{\iota \in \mathbf{I}} \mathbf{X}_{\iota}\) such that \(g(\iota) \in \mathbf{X}_{\iota}\) for all \(\iota \in \mathbf{J}\) , then there exists an extension \(f\) of \(g\) to \(\mathbf{I}\) such that \(f(\iota) \in \mathbf{X}_{\iota}\) for all \(\iota \in \mathbf{I}\) .

For each \(\iota\in I-J\) let \(T_{\iota}\) denote the term \(\tau_{y}(y\in X_{\iota})\) . Since \(X_{\iota}\neq\emptyset\) by hypothesis, we have \(T_{\iota}\in X_{\iota}\) for all \(\iota\in I-J\) (Chapter I, §4, no. 1). If G is the graph of \(g\) , the graph \(G\cup\left(\bigcup_{\iota\in I-J}\{(\iota,T_{\iota})\}\right)\) is the graph of a function which has the required properties, as is immediately verified.

COROLLARY 1. Let \((\mathbf{X}_{\iota})_{\iota \in \mathbf{I}}\) be a family of sets such that for each \(\iota \in \mathbf{I}\) we have \(\mathbf{X}_{\iota} \neq \emptyset\) . Then for each \(\alpha \in \mathbf{I}\) the projection \(\operatorname{pr}_{\alpha}\) is a mapping of \(\prod_{\iota \in \mathbf{I}} \mathbf{X}_{\iota}\) onto \(\mathbf{X}_{\alpha}\) .

Apply Proposition 5 to the subset \(J = \{\alpha\}\) of I and note that \(\mathbf{pr}_{\alpha}\) is the composition of the canonical mapping of \(X_{\alpha}^{\{\alpha\}}\) onto \(X_{\alpha}\) and the mapping \(\mathbf{pr}_{\{\alpha\}}\) .

COROLLARY 2. Let \((\mathbf{X}_t)_{i \in \mathbf{I}}\) be a family of sets. Then \(\prod_{i \in \mathbf{I}} \mathbf{X}_t = \varnothing\) if and only if there exists \(i \in \mathbf{I}\) such that \(\mathbf{X}_t = \varnothing\) .

If we have \(X_{t} \neq \varnothing\) for each \(t \in I\) , then it follows from Corollary 1 that \(\prod_{i \in I} X_{t} \neq \varnothing\) ; conversely, if \(\prod_{i \in I} X_{t} \neq \varnothing\) , the relation \(\Pr_{\alpha}\left(\prod_{i \in I} X_{t}\right) \subset X_{\alpha}\) shows that \(X_{\alpha} \neq \varnothing\) for each \(\alpha \in I\) .

Hence, if we have a family \((\mathbf{X}_{t})_{t\in\mathbf{I}}\) of non-empty sets, we may introduce (as an auxiliary constant) a function f with domain I such that \(f(t)\in\mathbf{X}_{t}\) for all \(t\in I\) . In practice one says: ``take an element \(x_{t}\) in each \(X_{t}\) ''. Intuitively, we have thus ``chosen'' an element \(x_{t}\) in each set \(X_{t}\) ; the introduction of the logical sign \(\tau\) and the criteria governing its use absolve us from the necessity of formulating an ``axiom of choice'' to legalize this operation (cf.~Summary of Results, §4, no. 10).

COROLLARY 3. Let \((\mathbf{X}_t)_{i\in I}\) and \((\mathbf{Y}_t)_{i\in I}\) be two families of sets having the same index set I. If \(\mathbf{X}_t\subset \mathbf{Y}_t\) for each \(i\in I\) , then

\[
\prod_ {i \in I} X _ {i} \subset \prod_ {i \in I} Y _ {i}.
\]

Conversely, if \(\prod_{i\in I}X_i\subset \prod_{i\in I}Y_i\) , and if \(X_{i}\neq \emptyset\) for each \(i\in I\) , then \(X_{i}\subset Y_{i}\) for each \(i\in I\) .

The first assertion is obvious, and the second follows from Proposition 1, Corollary 5, because we have then, for each \(\alpha \in I\) ,

\[
\mathrm{X} _ {\alpha} = \operatorname{pr} _ {\alpha} \left(\prod_ {i \in \mathrm{I}} \mathrm{X} _ {i}\right) \subset \operatorname{pr} _ {\alpha} \left(\prod_ {i \in \mathrm{I}} \mathrm{Y} _ {i}\right) = \mathrm{Y} _ {\alpha}.
\]

